\chapter{The Assignment}

% ============================================================
%  What the homework asks and how we are going to solve it
% ============================================================

\section{What is asked}

The homework asks us to build a small \textbf{JavaScript/HTML/CSS} program that shows how, in Ethereum, a \textit{private key} is turned into a \textit{public key}. The mechanism is \textbf{elliptic-curve scalar multiplication}:
\[
  Q = k \cdot G
\]
where $k$ is the private key (an integer), $G$ is a fixed public point of the curve (the \textit{base point}) and $Q$ is the public key (another point of the curve).

Ethereum uses the curve \texttt{secp256k1}, whose numbers are 256 bits long. Since the goal is to \underline{understand} the idea and not to build a real cryptosystem, we work on a \textbf{toy curve} with tiny numbers, so that every step can be checked by hand:
\[
  E:\; y^2 \equiv x^3 + 7 \pmod{17}, \qquad G = (1,5), \qquad k = 6.
\]
Note that the toy curve has the same equation as \texttt{secp256k1} ($y^2 = x^3 + 7$): only the modulus changes.

\begin{framedobs}{Why it is a ``one-way function''}
  The map $k \mapsto kG$ is
  \begin{itemize}
    \item \textbf{deterministic}: the same $k$ always gives the same $Q$;
    \item \textbf{easy forward}: computing $kG$ takes about $\log_2 k$ point operations;
    \item \textbf{hard backward}: given $Q$ and $G$, finding $k$ is the \textit{Elliptic Curve Discrete Logarithm Problem} (ECDLP), for which no efficient algorithm is known.
  \end{itemize}
  This is why it ``looks like'' a hash function. However, it has an algebraic structure (it is a group homomorphism) that a hash function must \underline{not} have: this difference is the core of the statistical part of the homework.
\end{framedobs}

\section{Deliverables, step by step}

\begin{enumerate}
  \item \textbf{Define the toy curve and the base point}: $p = 17$, $a = 0$, $b = 7$, $G = (1,5)$. Check that $G$ is really on the curve: $5^2 = 25 \equiv 8$ and $1^3 + 7 = 8$. \checkmark
  \item \textbf{Define the private key}: $k = 6$.
  \item \textbf{Implement point addition}: the group law $P + Q$ on the curve, including the special cases (doubling $P + P$, the point at infinity $\mathcal{O}$, $P + (-P) = \mathcal{O}$). It needs modular arithmetic and the modular inverse.
  \item \textbf{Compute and print} $G, 2G, 3G, \ldots, kG$ by repeated addition.
  \item \textbf{Faster method} for large $k$: the \textit{double-and-add} algorithm (the elliptic-curve version of square-and-multiply).
  \item \textbf{Print} $kG$ as the \textbf{TOY PUBLIC KEY}. With our data: $6G = (5,8)$.
  \item \textbf{Optional (guru level)}: a slowed-down animation of the algorithm on a canvas/SVG, showing the points of the curve on the $17 \times 17$ grid and each addition step.
  \item \textbf{Statistical exploration}: repeat the experiment with many values of $k$ (and other base points), study the distribution of the coordinates, of the bits, uniformity, input/output correlation and avalanche effect, and discuss which properties are meaningful for a one-way function and which ones belong to hash functions.
  \item \textbf{Publish} the result on the GitHub Pages site (\texttt{hw1/} folder of the \textit{Statistics Homework} site).
\end{enumerate}

\section{Roadmap of these notes}

\begin{itemize}
  \item \Cref{ch:background} gives the theory: modular arithmetic, elliptic curves over finite fields, the group law, scalar multiplication, ECDLP and how Ethereum uses all of this.
  \item \Cref{ch:worked} solves the toy example completely by hand, step by step.
  \item \Cref{ch:stats} explains what to measure in the statistical part, what we expect to see, and how to interpret it.
\end{itemize}
